\section{Case Study}
\label{sec:case-study}

The purpose of this section is to establish one claim: on the dynamics-encoding fragment, checking a temporal property set is a usable substitute for constructing a homomorphism, at a scale where constructing one by hand is unattractive. The section builds up to a system whose functional intent is arbitrary---a Turing machine, presented as a coupling of two zones through declared ports---then varies how its zones are built, and then hands the resulting conformance questions to a solver.

The order is deliberate:
\begin{enumerate}
    \item \textbf{Kit zones} (Section~\ref{sec:case-basics}). Four small reference systems---an incrementing counter, a saturating subtracter, a clear/add accumulator, and a zero-test---each presented as a five-tuple with a compiled fragment $\Phi_{\mathrm{dyn}}$ and a direct buildable. These fix notation and show what a compiled fragment looks like when it is short enough to read.
    \item \textbf{Alternate buildables} (Section~\ref{sec:case-alternates}). The same references held fixed while buildable encodings change, contrasting re-checking a fixed fragment against reconstructing a homomorphism $h=(h_S,h_I,h_O)$ by hand.
    \item \textbf{A coupled machine with arbitrary intent} (Section~\ref{sec:case-machine}). Two zones with ports---a tape and a finite control---a connectivity, and a resultant, related to a monolithic reference. Universality is definitional for this construction, so the demonstration covers any computable functional intent rather than a convenient sample. Zone variants, accepted and rejected, follow; zone conformance then lifts to the machine without re-examining the machine.
    \item \textbf{Solver-decided conformance} (Section~\ref{sec:case-solver}). The conformance question for each of those builds, encoded and decided mechanically: a conformance map returned when one exists, a proof of non-existence when it does not.
    \item \textbf{Matched ticks} (Section~\ref{sec:case-ticks}). What the fragment certifies and what it cannot, using Fibonacci as a sketch and a two-line negative result about temporal granularity.
\end{enumerate}

For readers less familiar with T3SD: a \emph{reference system} $Z_{\text{spec}}$ is a complete discrete-system table---states, inputs, outputs, next-state, readout---expressing required behavior. A \emph{buildable system} $Z_{\text{impl}}$ is another such table, often with richer internal state, intended to realize it. Classical checking asks for a homomorphism from $Z_{\text{impl}}$ onto $Z_{\text{spec}}$; here we compile $\Phi_{\mathrm{dyn}}$ from $Z_{\text{spec}}$ and check that property set instead. A \emph{coupling recipe} adds to a vector of such systems a connectivity---a one-to-one pairing of output ports with input ports---whose \emph{resultant} is again a single discrete system, with the unconnected ports as its interface.

Two scope notes. Every conformance and rejection claim in Sections~\ref{sec:case-machine} and~\ref{sec:case-ticks} is machine-checked in a proof assistant, including the negative ones; a rejection is a proof that no maps exist, not a failed search. The solver results of Section~\ref{sec:case-solver} are separate evidence: they concern a finite, tape-bounded instance, they are cross-checked against those theorems wherever one exists, and they are feasibility evidence rather than a scalability claim.

\subsection{Encoding IOR as Temporal Clauses}
\label{sec:case-basics}

The characteristic construction of Section~\ref{sec:trace-observability} applies to every IOR. For structured relations, it reduces to compact laws that look familiar to classical LTL:
\begin{align}
 \text{echo:}\quad
   &G\bigl(y(t)=f(t)\bigr), \label{eq:echo-phi}\\
 \text{one-tick delay:}\quad
   &y(0)=b_0\ \wedge\
     G\bigl(y(t{+}1)=f(t)\bigr), \label{eq:delay-phi}\\
 \text{run length:}\quad
   &y(t)=\phi_{\mathrm{out}}(f,t). \label{eq:rle-phi}
\end{align}
Each equation abbreviates an eligible-output relation $ER(f)$.
Note that equation~\eqref{eq:echo-phi} is memoryless. A Moore system with a fixed initial state cannot make $y(0)$ depend on an input first observed at tick zero, so its functionality cotyledon may be empty. Equation~\eqref{eq:delay-phi} adds the minimum delay required by state transitions and is always realizable by a one-bit register. The fragment therefore exposes the timing constraint already present in Wymore's trajectory semantics.

For Equation~\eqref{eq:rle-phi}, $\phi_{\mathrm{out}}(f,0)=(0,0)$;
tick $2k+1$ emits the bit and run length for symbol $f(2k)$; tick $2k+2$
holds the pair. The next subsection gives the run-length recurrence and the
two-tick pipeline shift. Finite as well as nondeterministic requirements use
the same construction with a restricted time set or a multi-valued $ER(f)$.

\subsection{Alternate Buildables: Same Specs, Changing Encodings}
\label{sec:case-alternates}

The preceding directs treat each reference in isolation with an identity buildable. The value of a compiled $\Phi_{\mathrm{dyn}}$ appears when the \emph{same} reference is held fixed while buildable structure evolves. Three alternates illustrate the pattern; the binary-shift counter is the structural-divergence exhibit.

\paragraph{Binary-shift counter for $Z_{\mathrm{inc}}$.}
\label{sec:case-shift-counter}
State is a little-endian bit list; enable $1$ performs binary increment with carry; enable $0$ holds; readout is the Nat value of the bits. There is no explicit count component---only the bit window.

\begingroup
\renewcommand{\arraystretch}{1.2}
\footnotesize
\begin{center}
\begin{tabularx}{0.92\linewidth}{@{}>{\raggedright\arraybackslash}p{2.8cm}Y@{}}
    \rowcolor{tableheader}
    \headingfont\bfseries Component & \headingfont\bfseries $Z_{\mathrm{inc}}^{\mathrm{shift}}$ \\
    \addlinespace[2pt]
    $S$ & finite bit lists (little-endian) \\
    \addlinespace[2pt]
    $I,O$ & $\{0,1\}$, $\mathbb{N}$ \\
    \addlinespace[2pt]
    $\delta(\mathit{bits},1)$ & binary increment with carry \\
    \addlinespace[2pt]
    $\delta(\mathit{bits},0)$ & $\mathit{bits}$ \\
    \addlinespace[2pt]
    $\rho(\mathit{bits})$ & Nat value of $\mathit{bits}$ \\
    \addlinespace[2pt]
    Maps to $Z_{\mathrm{inc}}$ & $h_S=\mathrm{value}$,\ $h_I=\mathrm{id}$,\ $h_O=\mathrm{id}$ \\
\end{tabularx}
\end{center}
\captionof{table}{Binary-shift counter: no explicit count field; $h_S$ interprets the bit list.}
\label{tab:z-inc-shift}
\endgroup
\vspace{4pt}

\paragraph{Workflow contrast.}
\textit{Assertional:} $\Phi_{Z_{\mathrm{inc}}}$ is unchanged from the direct buildable; verifying the shift encoding means re-checking that same fragment (equivalently, exhibiting/verifying a Def.~4.3 witness on finite abstractions, or checking the exhibited maps on $\mathbb{N}$).
\textit{Classical:} $h_S$ is the bit-list valuation---not a component projection. Transition preservation requires proving that binary increment raises the valuation by one; each structural change to the bit encoding forces that argument to be revisited.

\paragraph{Instrumented decrement for $Z_{\mathrm{dec}}$.}
State $(n,k)$ tracks the current value and a pulse count of enable-$1$ steps; readout is $n$. Classical $h_S$ is the projection $(n,k)\mapsto n$. Assertional checking reuses $\Phi_{Z_{\mathrm{dec}}}$; the extra counter is invisible to the fragment so long as the projection is preserved.

\paragraph{Two zero-test encodings for $Z_{\mathrm{jz}}$.}
An elaborated latch stores $(\mathit{flag},\mathit{lastSample})$ and projects to the flag; a dual encoding stores the flag twice and projects to the first copy. Both are checked against the same $\Phi_{Z_{\mathrm{jz}}}$. Classical witnesses remain projections; the point is that instrumentation does not force a new property set.

\paragraph{What checking means.}
Across these alternates, $\mathit{specState}$ atoms are interpreted relative to the chosen $h_S$. Revising an elaboration does not change the compiled fragment indexed by $Z_{\text{spec}}$; one re-checks the same property set. The binary-shift counter is the case where that re-check is paired with a non-trivial classical $h_S$.

\subsection{A Coupled Machine with Arbitrary Functional Intent}
\label{sec:case-machine}

The kit zones of Section~\ref{sec:case-basics} are single systems whose intent is stated in one line. Engineering intent is rarely like that, and a reader is entitled to ask whether the compiled fragment says anything once intent becomes demanding. This section answers that in the strongest available form: it builds a system whose functional intent is \emph{any} computable behavior, presents it as a coupling of two zones through declared ports, and shows that conformance of the coupling to a monolithic reference is exactly satisfaction of the compiled fragment.

The system is a Turing machine. Nothing about the presentation is computational in flavor---no program, no software, no domain story. There are two zones with ports, a wiring table, and a resultant, in the ordinary Wymorean sense. What the choice buys is that universality comes for free: a tape zone coupled to a finite control zone realizes any computable functional intent, so no reader can object that the demonstration was chosen to be easy.

\subsubsection{Two zones, stated as discrete systems with ports}
\label{sec:case-machine-zones}

Fix a finite symbol set $\Gamma$ with a blank, and a finite label set $\Lambda$. A \emph{table} $M$ assigns to each label and scanned symbol either a halt or a triple (next label, symbol to write, head move).

\paragraph{Tape zone $Z_{\mathrm{tape}}$.}
State is a tape: a finite left part, a scanned symbol, a finite right part, extended with blanks as the head runs off either end. The zone has one input port carrying a command---either \emph{hold} or \emph{act}$(w,m)$, meaning write $w$ then move $m$---and two output ports, both carrying the scanned symbol. Two output ports are not redundancy: coupling connectivity is one-to-one, so exposing the scanned symbol to the control \emph{and} at the system boundary requires distinct ports.

\begingroup
\renewcommand{\arraystretch}{1.25}
\footnotesize
\begin{center}
\begin{tabularx}{0.95\linewidth}{@{}>{\raggedright\arraybackslash}p{3.0cm}Y@{}}
    \rowcolor{tableheader}
    \headingfont\bfseries Component & \headingfont\bfseries $Z_{\mathrm{tape}}$ \\
    \addlinespace[2pt]
    $S$ & tapes $(\ell, a, r)$ over $\Gamma$ \\
    \addlinespace[2pt]
    Input ports & $\mathit{cmd}$ \\
    \addlinespace[2pt]
    Port value sets & $I_{\mathit{cmd}} = \{\mathit{hold}\} \cup \{\mathit{act}(w,m) \mid w\in\Gamma,\ m\in\{L,R,S\}\}$ \\
    \addlinespace[2pt]
    Output ports & $\mathit{scan}$, $\mathit{window}$ \\
    \addlinespace[2pt]
    Port value sets & $O_{\mathit{scan}} = O_{\mathit{window}} = \Gamma$ \\
    \addlinespace[2pt]
    $\delta(t, \mathit{cmd})$ & apply the command: write, then move, extending with blanks \\
    \addlinespace[2pt]
    $\delta(t, \mathrm{none})$ & $t$ \\
    \addlinespace[2pt]
    $\rho(t)$ & scanned symbol on both output ports \\
\end{tabularx}
\end{center}
\captionof{table}{Tape zone $Z_{\mathrm{tape}}$: one command port in, the scanned symbol on two ports out.}
\label{tab:z-tape}
\endgroup
\vspace{4pt}

\paragraph{Control zone $Z_{\mathrm{ctl}}$.}
Readout in a discrete system is a function of state alone. In a feedback coupling this has a consequence that is easy to miss and impossible to work around: the control cannot both observe the scanned symbol and emit a command derived from it within one tick, because the command it offers this tick was fixed by the state it was already in. The control is therefore a two-phase system. In a \emph{sense} phase it latches the action the table selects for the symbol it is being shown; in an \emph{act} phase it emits that action on its command port and returns to sensing. A halted state emits \emph{hold} forever. One machine step consequently occupies two ticks of the coupling, a point we return to in Section~\ref{sec:case-ticks}.

\begingroup
\renewcommand{\arraystretch}{1.25}
\footnotesize
\begin{center}
\begin{tabularx}{0.95\linewidth}{@{}>{\raggedright\arraybackslash}p{3.0cm}Y@{}}
    \rowcolor{tableheader}
    \headingfont\bfseries Component & \headingfont\bfseries $Z_{\mathrm{ctl}}$ \\
    \addlinespace[2pt]
    $S$ & $\{\mathit{sense}(q)\} \cup \{\mathit{act}(q,c)\} \cup \{\mathit{halted}\}$, $q\in\Lambda$ \\
    \addlinespace[2pt]
    Input ports & $\mathit{sym}$, $\mathit{load}$ \\
    \addlinespace[2pt]
    Port value sets & $I_{\mathit{sym}} = \Gamma$, \quad $I_{\mathit{load}} = \Lambda \cup \{\mathrm{none}\}$ \\
    \addlinespace[2pt]
    Output ports & $\mathit{cmd}$, $\mathit{halted}$ \\
    \addlinespace[2pt]
    Port value sets & $O_{\mathit{cmd}} = I_{\mathit{cmd}}$ of Table~\ref{tab:z-tape}, \quad $O_{\mathit{halted}} = \{0,1\}$ \\
    \addlinespace[2pt]
    $\delta(\mathit{sense}(q), (a, \mathrm{none}))$ & $\mathit{act}(q', \mathit{act}(w,m))$ if $M(q,a)=(q',w,m)$; else $\mathit{halted}$ \\
    \addlinespace[2pt]
    $\delta(\mathit{act}(q,c), (a, \mathrm{none}))$ & $\mathit{sense}(q)$ \\
    \addlinespace[2pt]
    $\delta(s, (a, q_0))$ & $\mathit{sense}(q_0)$ \quad (an external load overrides the table) \\
    \addlinespace[2pt]
    $\rho(s)$ & the command $s$ offers, and whether $s$ is halted \\
\end{tabularx}
\end{center}
\captionof{table}{Control zone $Z_{\mathrm{ctl}}$: sense latches the table's decision, act emits it.}
\label{tab:z-ctl}
\endgroup
\vspace{4pt}

\subsubsection{The coupling recipe and its resultant}
\label{sec:case-machine-coupling}

The two zones form a connectable vector $(Z_{\mathrm{tape}}, Z_{\mathrm{ctl}})$. A coupling recipe adds a connectivity: a one-to-one correspondence between output ports and input ports that is proper in both directions, so that some ports on each side remain free and constitute the boundary of the resultant. Here the connectivity is the two-wire table below.

\begingroup
\renewcommand{\arraystretch}{1.25}
\footnotesize
\begin{center}
\begin{tabularx}{0.95\linewidth}{@{}>{\raggedright\arraybackslash}p{4.4cm}>{\raggedright\arraybackslash}p{4.4cm}Y@{}}
    \rowcolor{tableheader}
    \headingfont\bfseries From output port & \headingfont\bfseries To input port & \headingfont\bfseries Carries \\
    \addlinespace[2pt]
    $Z_{\mathrm{ctl}}.\mathit{cmd}$ & $Z_{\mathrm{tape}}.\mathit{cmd}$ & tape command \\
    \addlinespace[2pt]
    $Z_{\mathrm{tape}}.\mathit{scan}$ & $Z_{\mathrm{ctl}}.\mathit{sym}$ & scanned symbol \\
\end{tabularx}
\end{center}
\captionof{table}{Connectivity of the recipe: a pure feedback coupling of the two zones.}
\label{tab:machine-wires}
\endgroup
\vspace{4pt}

Three ports are left unconnected and are therefore the interface of the coupled system: the control's $\mathit{load}$ input, the control's $\mathit{halted}$ output, and the tape's $\mathit{window}$ output. The resultant $Z_{\mathrm{mach}} = \mathrm{RSY}$ of the recipe is again a discrete system in the ordinary sense: its state is a pair (tape, control phase); its input is whatever arrives on $\mathit{load}$; its output is the pair (scanned symbol, halted flag) offered on the two free output ports. Its next-state function is the componentwise next state of the two zones, each fed by the values resolved on its input ports---which for a connected port is the value the driving output port currently offers, and for a free port is the value supplied by the environment.

\begin{figure}[!ht]
  \centering
  \begin{tikzpicture}[
    >=Stealth,
    font=\footnotesize,
    zone/.style={draw=blue!65!black, fill=blue!6, rounded corners=2pt,
      align=center, minimum height=1.35cm, minimum width=3.1cm, thick},
    wire/.style={->, thick, draw=black!70},
    bnd/.style={->, thick, draw=black!45, dashed},
    port/.style={font=\scriptsize, align=center}
  ]
    \node[zone] (tape) {$Z_{\mathrm{tape}}$\\[1pt]{\scriptsize state: tape}};
    \node[zone, right=3.4cm of tape] (ctl) {$Z_{\mathrm{ctl}}$\\[1pt]{\scriptsize state: sense / act / halted}};

    \draw[wire] ([yshift=9pt]ctl.west) -- node[above, port] {$\mathit{cmd}$} ([yshift=9pt]tape.east);
    \draw[wire] ([yshift=-9pt]tape.east) -- node[below, port] {$\mathit{scan}$} ([yshift=-9pt]ctl.west);

    \node[port, left=1.0cm of tape] (win) {$\mathit{window}$};
    \draw[bnd] (tape.west) -- (win);
    \node[port, right=1.0cm of ctl] (hal) {$\mathit{halted}$};
    \draw[bnd] (ctl.east) -- (hal);
    \node[port, above=0.75cm of ctl] (load) {$\mathit{load}$};
    \draw[bnd] (load) -- (ctl.north);

    \begin{scope}[on background layer]
      \node[draw=black!35, dashed, rounded corners=3pt, inner sep=13pt,
        fit=(tape)(ctl)] (box) {};
    \end{scope}
    \node[below=2pt of box, font=\scriptsize] {resultant $Z_{\mathrm{mach}}$; dashed ports are the system boundary};
  \end{tikzpicture}
  \caption{The coupling recipe of Table~\ref{tab:machine-wires}. Solid arrows are the connectivity; dashed ports are unconnected and form the interface of the resultant.}
  \label{fig:machine-coupling}
\end{figure}

\subsubsection{The reference, and what conformance means here}
\label{sec:case-machine-reference}

The functional intent is stated separately, as a single reference system $Z_{\mathrm{ref}}$ over configurations, with no ports and no internal structure: state is a pair (tape, control phase), input is the load command, output is the pair (scanned symbol, halted flag), and one tick advances the tape by the command the control currently offers while stepping the control on the symbol currently under the head. This is intent as a systems engineer would write it---a table of required behavior---and it is deliberately not the coupling.

Compiling the fragment of Section~\ref{sec:dynamics-encoding-fragment} from $Z_{\mathrm{ref}}$ yields one open-readout clause per configuration, one guided-step clause per configuration and load value, and one autonomous clause per configuration. The bi-implication then says that the coupled machine satisfies that property set exactly when it realizes the reference. In the machine-checked development the coupling does satisfy it, and in fact the relation is tighter than required: the state map that reads off the two component states, together with the port readings on the boundary, is a bijection, so the coupling and the reference are isomorphic. Nothing was lost or added by decomposing intent into two zones and a wiring table.

Two remarks matter for how much this example proves.

\emph{Universality is definitional, not asserted.} The construction is parametric in the table $M$. A tape coupled to a finite control realizes any computable functional intent by the standard argument; the case study therefore covers arbitrary computable behavior rather than a sample of it. This is the property the earlier kit examples cannot claim, and it is why the machine is the spine of the case study.

\emph{The coupling is not a re-skin of the reference.} The reference is a flat table over configurations; the coupling is two independently defined zones whose only interaction is through the two declared wires, with the tape unable to observe the control's label and the control unable to observe anything but the scanned symbol. That the resultant of the recipe agrees tick-for-tick with the flat table is a theorem about the wiring, and it is the theorem the compiled fragment certifies.

\subsubsection{The same intent, built differently}
\label{sec:case-machine-variants}

A single verified design demonstrates little about a verification method. The engineering question is which of several candidate builds conform, so we vary the zones and keep the reference fixed. Each variant below is a legitimate design decision or a plausible defect, and each verdict is a theorem.

\paragraph{Instrumented tape (accepted).}
A tape that additionally maintains an actuation count, exposed on no port. The zone-level state map forgets the count. It is a homomorphism and not an isomorphism: two implementation states differing only in the count are identified, which is the ordinary situation when a build carries internal structure its reference does not mention.

\paragraph{Re-encoded control (accepted).}
The same control under a different state encoding---a tagged sum rather than a record---with the dynamics transported along the encoding. The zone-level relation is an isomorphism. This variant exists to make a specific point: the compiled fragment cannot see state encodings at all, for reasons developed in Section~\ref{sec:invariance}.

\paragraph{Frozen tape (rejected).}
A tape that accepts commands on its port and never actuates them---a design that forgot the actuator. No state, input and output map whatever makes it a realization of the tape zone. The proof is short and worth stating, because it shows what a rejection looks like: pick an implementation state that the state map sends to a tape whose scanned symbol is $a$, and an implementation input the input map sends to the command \emph{write $b$, stay}. The frozen tape does not move, so the implementation state is unchanged; the reference must now read $b$ where it reads $a$. No choice of maps can absorb that.

\paragraph{Blind tape (rejected).}
A tape that actuates correctly but always reports a blank on its window. Its readout is constant, so it cannot cover two reference states that read out differently, and the state map is required to cover all of them.

\paragraph{Over-reporting build against a silent reference (rejected).}
The clause family that is easiest to overlook is the closed readout, so it is worth one variant of its own. Consider a reference tape whose window is undefined until the head has moved off the left end---a sensor that is not yet valid, stated at the reference level as no output. A build that reports on every tick fails, because a reference state with no output can only be covered by an implementation state that also produces none. Conformance constrains silence as well as speech.

\subsubsection{Verify the zones, obtain the machine}
\label{sec:case-machine-lift}

The last movement is the one that matters for scale. Swap \emph{both} accepted variants into the recipe of Table~\ref{tab:machine-wires} at once, leaving the connectivity untouched. Because each zone of the rebuilt recipe is a port-preserving homomorphic image of the corresponding zone of the original, the resultant of the rebuilt recipe is a port-preserving homomorphic image of the original resultant, and the boundary is unchanged: the same free ports, carrying the same values. Composing that with the realization of Section~\ref{sec:case-machine-reference} gives conformance of the rebuilt machine to the reference---and the machine-level argument is never re-opened. The reasoning is entirely at the zones and at the wiring.

This is the property that makes zone-level checking worth automating. A solver that can only settle small systems is still useful if zone verdicts assemble into system verdicts, and Section~\ref{sec:case-solver} takes that seriously by handing the zone-level questions to a solver.

\subsection{Finite Checking}
\label{sec:case-solver}

The unbounded machine is established symbolically. Finite projections provide
automated checks of bounded behavior. Capping the count at $n$ preserves the
same transition whenever the current count is below $n$; the trace $0,0,1$
remains below a cap of four.

For any finite projection, satisfaction of $\Phi_{\mathrm{dyn}}$ becomes a
propositional search for the three surjective maps subject to the readout,
guided-step, autonomous-step, and surjectivity clauses. A satisfying assignment
is a Wymore homomorphism; unsatisfiability establishes that no such map exists.
The matched- and mismatched-tick instances in Section~\ref{sec:case-ticks}
exercise both outcomes. The same encoding can be generated for each finite
component or abstraction in a larger coupling.

\subsection{Matched Ticks: What the Fragment Can and Cannot Say}
\label{sec:case-ticks}

The machine of Section~\ref{sec:case-machine} invites a natural question, and answering it exposes the sharpest limitation of the fragment. The question: since the construction is parametric in the table, there is a table for which the machine computes Fibonacci numbers---so can the compiled fragment certify \emph{that}?

\subsubsection{Fibonacci, as a sketch}
\label{sec:case-fib}

The instinct to reach for Fibonacci is worth following only as far as it goes, so we state the sketch and stop. Take the table $M_{\mathrm{fib}}$ of a standard machine that, started with $n$ in unary on the tape, halts with the $n$th Fibonacci number written in unary. Instantiating Section~\ref{sec:case-machine} at $M_{\mathrm{fib}}$ gives a coupled system, a monolithic reference, and---by the theorems already established---a compiled fragment that the coupling satisfies exactly when it realizes that reference. In that sense there is nothing left to do: a Fibonacci machine is a machine, and the case study covered arbitrary tables on purpose.

But it is worth being precise about what such a certificate would and would not say, because the gap is instructive rather than embarrassing. Conformance to the reference says: \emph{this build takes the same steps as that table}. It does not say \emph{this build eventually halts with $F_n$ on the tape}. The latter is a liveness statement about a run, and the fragment has no $F$ operator; it has readout clauses and next-state clauses and nothing else. One can recover a bounded version---after a stated number of ticks the window reads $F_n$---by unrolling the reference, and that is a genuine and checkable property. What one cannot do is state ``eventually done'' inside the fragment at all. The honest summary is that $\Phi_{\mathrm{dyn}}$ certifies \emph{step-for-step correctness against intent}, and that termination and other liveness intent must be argued in a richer dialect, outside the bi-implication. Section~\ref{sec:expressivity} takes up the trade-off.

\subsubsection{Steps are not ticks, and the fragment counts ticks}
\label{sec:case-ticks-witness}

There is a second, less obvious lesson available from the same place. Recall from Section~\ref{sec:case-machine-zones} that one machine step occupies two ticks of the coupling, because the control must sense before it can act. That was forced by readout depending on state alone. The reference of Section~\ref{sec:case-machine-reference} was therefore written at the coupling's granularity---two ticks per machine step---and not at the granularity a systems engineer would naturally choose when writing down intent, which is one row per machine step.

That choice was not cosmetic. A Wymore homomorphism relates ticks to ticks: the step law equates one implementation transition with one reference transition. It has no room for a build that takes two ticks to accomplish what the reference does in one, however obviously correct that build is. This is worth exhibiting on the smallest possible pair, so that the obstruction is visible rather than asserted.

Let the reference be a one-bit register that inverts on every driven tick and reports its bit. Let the build be the same register with a sense/act phase, inverting every \emph{second} driven tick and reporting its bit throughout---exactly the two-phase pattern the control zone uses. There is no conformance map, and the argument is three lines. The readout law forces the state map to factor through the bit: any two build states with the same bit are sent to the same reference state, so the map cannot distinguish the sense phase from the act phase. But the step law applied to the sense state requires its successor---the act state, with the same bit---to be sent to the inverted reference state. So one reference state must equal its own inversion.

The two-phase build is not wrong; the reference was written at the wrong granularity. Restate the same intent at two ticks per logical inversion and the identity maps are a conformance map, indeed an isomorphism. Both verdicts are machine-checked, and both appear in Table~\ref{tab:solver-results} as the pair \texttt{flip-one-tick} and \texttt{flip-matched-ticks}: the same build, rejected against one reference and accepted against the other.

\subsubsection{What to do about it}
\label{sec:case-ticks-practice}

The practical reading is a modeling obligation, and it should be stated as such rather than buried. \emph{Reference and build must agree on what a tick is.} Two consequences for practice follow.

First, when intent is written at a coarser granularity than the build can achieve---the usual case, since intent is written before the technology that will realize it is chosen---the reference must be expanded to the build's granularity before conformance is even the right question. The expansion is mechanical and it is where the design's internal phasing becomes visible in the specification. For the machine, this is why the reference has a phase component at all: not because intent cares about phases, but because the technology needs two ticks and lock-step conformance forces intent to say so.

Second, this is a real limitation of the classical relation rather than of the temporal encoding, and the bi-implication is what makes that clear. Because satisfaction of the fragment and existence of a homomorphism are the same condition, the fragment inherits exactly the lock-step character of the classical homomorphism---no more and no less. A verification method that tolerated temporal abstraction would relate a build's $k$ ticks to a reference's one, which is a weaker relation than Wymore's and would need its own theory. We do not develop one here; we note that the fragment's rejection of the two-phase build is faithful reporting of the classical relation, and that a reader who wants tolerance for differing rates should read our negative result as identifying precisely what would have to be relaxed.

